\documentclass[11pt]{article}
\usepackage{amsmath,amssymb,amsthm,mathrsfs}
\usepackage[margin=1in]{geometry}
\usepackage{hyperref}

\newtheorem{theorem}{Theorem}
\newtheorem{proposition}[theorem]{Proposition}
\newtheorem{lemma}[theorem]{Lemma}
\newtheorem{corollary}[theorem]{Corollary}
\theoremstyle{definition}
\newtheorem{remark}[theorem]{Remark}
\newtheorem{convention}[theorem]{Convention}
\newtheorem{definition}[theorem]{Definition}

\DeclareMathOperator{\ord}{ord}
\DeclareMathOperator{\ad}{ad}
\newcommand{\hgt}{\operatorname{ht}}   % NOT \ht -- collides with a TeX primitive
\newcommand{\Z}{\mathbb{Z}}
\newcommand{\Q}{\mathbb{Q}}
\newcommand{\Lam}{\Lambda}
\newcommand{\F}{\mathcal{F}}
\newcommand{\D}{\mathcal{D}}

\title{The order of $R_e(t)$ over the sublattice subring is still infinite\\
{\large and the two-parameter commutator $[R_e(t),R_f(s)]$}}
\author{Clio}
\date{7 September 2026}

\begin{document}
\maketitle

\begin{center}\small
\fbox{\parbox{0.92\textwidth}{\centering
\textbf{Author:} Clio \quad\textbullet\quad \textbf{Date:} 7 September 2026 \quad\textbullet\quad
\textbf{Recipient:} Rick (cc Robin)\\
\textbf{Title:} The order of $R_e(t)$ over the sublattice subring is still infinite,
and the two-parameter commutator $[R_e(t),R_f(s)]$\\
\textbf{Corresponds to:} \texttt{clio-vega/proofs@089bffd} (paper),
\texttt{@143197f} (registry nodes)\\[4pt]
\textbf{Annotation, 8 September 2026 (self-review, \texttt{clio-vega/rick-review@8223c88}, finding F2):}
Corollary~4.2(iv) --- and the sentence in the abstract above --- assert that
$[R_e(t),R_e(s)]$ has a nonzero two-bead part off $s=t$; the abstract's version also drops
the exclusion $ts=1$ that (iv) states.  \emph{Both are false for $e\in\{1,2\}$.}  The
two-assignment sum $t^{P-k}s^{Q}-t^{P}s^{Q+k}+t^{Q+k}s^{P}-t^{Q}s^{P-k}$ vanishes identically
\emph{iff} $k=0$ or $Q=P-k$, and $Q=P-k$ is \emph{forced} when $e=f\le2$: for $e=f=1$ the
interval $(c,c+1)$ is empty so $k=0$; for $e=f=2$, $k=\pm1$ forces $b=c\mp1$, whence
$P=\mathbf 1[c\in M]=1$ and $Q=\mathbf 1[c+1\in M]=0$ by legality.  Verified: $0$ nonzero
$e=f$ two-bead entries for $e\in\{1,2\}$ up to $|\lambda|\le8$, against $31$ for $e=3$ and
$223$ for $e=5$.  The closing clause of the proof of (iv), ``this is not identically zero'',
is the unwarranted step.  Correct statement: \emph{for $e\ge3$}, $[R_e(t),R_e(s)]$ has a
nonzero two-bead part off the two loci $s=t$ and $ts=1$; for $e\in\{1,2\}$ the two-bead part
is identically zero for all $t,s$.  This was visible in a number already printed in
\S\ref{sec:verif}: the count $61$ splits as $47\,(e{=}f{=}4)+14\,(e{=}f{=}3)+0+0$, and the
two zeros are the defect.  Theorem~4.1 itself is unaffected and was re-verified $1829/1829$
on an independent border-strip engine.
}}
\end{center}
\bigskip


\begin{abstract}
Let $R_e(t)$ add a connected $e$-ribbon with weight $t^{\hgt}$ to a Schur function, and let
$\Lam^{(e)}=\Q(t)[p_e,p_{2e},\dots]$.  Question Q96 asked whether the infinite differential
order of $R_e(t)$ --- proved over the full base $\Lam$ --- is an artifact of measuring
against \emph{all} the power sums, and collapses to a finite order when the base is shrunk
to the sublattice subring $\Lam^{(e)}$.

\textbf{The answer is no, and emphatically so.}  We prove the sharper statement that the
order is already infinite over the \emph{single-generator} subring $\Q(t)[p_e]$:
for every $e\ge2$, every $d\ge1$ and $\mu_d=\bigl((d+1)e-1,\,1\bigr)$,
\[
  \bigl\langle s_{\mu_d},\ \ad(M_{p_e})^{\,d}\bigl(R_e(t)\bigr)\,s_\varnothing\bigr\rangle
  \;=\;(-1)^{d}\,(1+t).
\]
Since $\Q(t)[p_e]\subseteq\Lam^{(e)}\subseteq\Lam$ and order is monotone \emph{decreasing}
in the base, this forces $\ord_{\Q(t)[p_e]}=\ord_{(e)}=\ord=\infty$ for every $t\neq-1$.
The deformation is genuinely transverse to the sublattice; the specialness of the anchor
$t=-1$ is not a choice of coordinates.

Three by-products.  (i) The witness is a \emph{single-channel} matrix element: for this
source and target there is exactly one legal sequence of $e$-moves, forced by a
flow-conservation argument on the abacus, and it is unique \emph{as an ordered sequence}.
The whole computation is therefore a $(d+1)$-term alternating sum with two distinct
weights, $t$ and $-1$.  (ii) The same witness re-proves the earlier full-base result
(\cite[Cor.~4.9]{Clio0906}) uniformly in the parameter, including the scalar $c=0$, which
previously needed a separate theorem.  (iii) We derive the two-parameter cross-rank
commutator $[R_e(t),R_f(s)]$ in closed form.  Its two-bead sector depends on the two
parameters only through the product $ts$, and vanishes identically exactly on the locus
$ts=1$; its one-bead sector is a genuinely \emph{two-coloured} dressing --- each half of
the interval carries the parameter of the operator that traversed it --- which is why it
does not collapse to a scalar multiple of a single dressed ribbon move when $s\neq t$.  We
also correct a case split: $[R_e(t),R_e(s)]$ has a \emph{nonzero} two-bead part when
$t\neq s$, so the ``$e=f\Rightarrow0$'' clause of \cite[Thm.~2.3(2)]{Clio0907} is a
one-parameter statement only.

The argument is elementary abacus combinatorics throughout.  In particular it does not
pass through Lam's LLT map $\Phi$ at any point (see Remark~\ref{rem:lam}), and no claim
about the literature is made anywhere in this note.

Verified by two code-disjoint engines; see \S\ref{sec:verif}.
\end{abstract}

\section{Setting}

\begin{convention}\label{conv:main}
$\Lam=\Q(t)[p_1,p_2,\dots]$ is the ring of symmetric functions over $\Q(t)$, with Schur
basis $\{s_\lambda\}$.  For a partition $\lambda$ the \emph{Maya set} is
$M(\lambda)=\{\lambda_j-j:j\ge1\}\subset\Z$; its elements are \emph{beads}.  We write
$|M\rangle$ for the corresponding $s_\lambda$ and identify $\Lam$ with the charge-$0$
semi-infinite wedge $\F$.  For $g\ge1$,
\[
  R_g(t)\,s_\lambda=\sum_{\mu}t^{\hgt(\mu/\lambda)}s_\mu,
\]
the sum over $\mu\supset\lambda$ with $\mu/\lambda$ a connected border strip of size $g$
and $\hgt=(\#\text{rows})-1$.  $M_f$ denotes multiplication by $f\in\Lam$.
\end{convention}

\begin{proposition}[abacus form; {\cite[Prop.~1.2]{Clio0907}}]\label{prop:abacus}
Call $(b,g)$ a \emph{legal $g$-move} on the Maya set $M$ if $b\in M$ and $b+g\notin M$,
and put $M\cdot(b,g)=(M\setminus\{b\})\cup\{b+g\}$.  Then
\[
  R_g(t)\,|M\rangle=\sum_{(b,g)\ \mathrm{legal}}t^{\,\#(M\cap(b,b+g))}\,\bigl|M\cdot(b,g)\bigr\rangle ,
\]
a finite sum; $b\mapsto M\cdot(b,g)$ is a bijection from legal $g$-moves onto the $\mu$
with $\mu/\lambda$ a connected $g$-ribbon, and $\hgt(\mu/\lambda)=\#(M\cap(b,b+g))$.  Here
$(b,b+g)$ is the \emph{open} interval.
\end{proposition}

\begin{proposition}[{Murnaghan--Nakayama; \cite[Q75]{Clio0903}}]\label{prop:mn}
$R_g(-1)=M_{p_g}$ for every $g\ge1$.
\end{proposition}

\noindent
This is the load-bearing input and it was re-verified at the start of the session against
an independent implementation of $p_g\cdot s_\lambda$ in $\Q[p_1,p_2,\dots]$
(\S\ref{sec:verif}, item 2).

\section{Order over a subring, and why the generators suffice}

\begin{definition}[order filtration over a base]\label{def:ord}
Let $B\subseteq\Lam$ be a $\Q(t)$-subalgebra containing $1$.  For a $\Q(t)$-linear
operator $X$ on $\Lam$ set $\D^{\le-1}_B=\{0\}$ and, for $d\ge0$,
\[
  \D^{\le d}_B=\bigl\{X:\ [X,M_f]\in\D^{\le d-1}_B\ \ \text{for every }f\in B\bigr\},
\]
equivalently $\bigl[\cdots[[X,M_{f_0}],M_{f_1}],\dots,M_{f_d}\bigr]=0$ for all
$f_0,\dots,f_d\in B$.  Put $\ord_B(X)=\min\{d:X\in\D^{\le d}_B\}$, and $\infty$ if there is
no such $d$.  We write $\ord=\ord_\Lam$ and $\ord_{(e)}=\ord_{\Lam^{(e)}}$ for
$\Lam^{(e)}=\Q(t)[p_e,p_{2e},p_{3e},\dots]$.
\end{definition}

$\Lam$ is free as a $\Lam^{(e)}$-module (a basis is given by the monomials in the $p_m$
with $e\nmid m$), so this is the Grothendieck order filtration of $X$ as a differential
operator on $\Lam$ \emph{over the base $B$}.  Domain and codomain are all of $\Lam$
throughout; we are changing the base, not restricting the domain.  (Restricting the domain
is not an option: $R_2(t)s_\varnothing=s_{(2)}+t\,s_{(1,1)}=\tfrac12(p_1^2+p_2)+\tfrac
t2(p_1^2-p_2)\notin\Lam^{(2)}$, so $R_e(t)$ does not preserve $\Lam^{(e)}$.)

\begin{lemma}[monotonicity in the base]\label{lem:mono}
If $B\subseteq B'$ are subalgebras of $\Lam$ then $\D^{\le d}_{B'}\subseteq\D^{\le d}_{B}$
for every $d$, hence $\ord_{B}(X)\le\ord_{B'}(X)$ for every $X$.  In particular
\[
  \ord_{\Q(t)[p_e]}(X)\ \le\ \ord_{(e)}(X)\ \le\ \ord(X).
\]
\end{lemma}

\begin{proof}
Induction on $d$.  For $d=-1$ both sides are $\{0\}$.  If $X\in\D^{\le d}_{B'}$ and $f\in
B\subseteq B'$ then $[X,M_f]\in\D^{\le d-1}_{B'}\subseteq\D^{\le d-1}_{B}$ by the inductive
hypothesis; so $X\in\D^{\le d}_{B}$.  The displayed chain is the case
$\Q(t)[p_e]\subseteq\Lam^{(e)}\subseteq\Lam$.
\end{proof}

So a lower bound proved over the \emph{smallest} base is the strongest of the three.  That
is what we will prove.

The next lemma is the reduction the plan asked to be established rather than imported.  It
is standard for a polynomial base, but it is load-bearing here, so we prove it.

\begin{lemma}[generators suffice]\label{lem:gen}
Let $B$ be generated as a $\Q(t)$-algebra with $1$ by a subset $S\subseteq\Lam$.  Define
$\mathcal E^{\le-1}=\{0\}$ and $\mathcal E^{\le d}=\{X:[X,M_g]\in\mathcal E^{\le d-1}\
\forall g\in S\}$.  Then $\mathcal E^{\le d}=\D^{\le d}_B$ for every $d\ge-1$.
Consequently $\ord_B(X)\le d$ if and only if every $(d+1)$-fold nested bracket of $X$
against multiplication operators $M_g$, $g\in S$, vanishes.
\end{lemma}

\begin{proof}
Each $\mathcal E^{\le c}$ is a $\Q(t)$-subspace, since a nested bracket is linear in $X$.

\emph{Step 1: $\mathcal E^{\le c}$ is stable under multiplication by any $M_h$, $h\in\Lam$,
on either side.}  Induct on $c$.  For $c=-1$ there is nothing to prove.  Let
$Y\in\mathcal E^{\le c}$ and $g\in S$.  Multiplication operators commute with each other,
so $[M_h,M_g]=0$ and
\[
  [\,YM_h,\ M_g\,]=[Y,M_g]\,M_h+Y\,[M_h,M_g]=[Y,M_g]\,M_h .
\]
By definition $[Y,M_g]\in\mathcal E^{\le c-1}$, so by the inductive hypothesis
$[Y,M_g]M_h\in\mathcal E^{\le c-1}$; hence $YM_h\in\mathcal E^{\le c}$.  The argument for
$M_hY$ is identical.

\emph{Step 2: the induction on $d$.}  $\mathcal E^{\le-1}=\D^{\le-1}_B$.  Assume
$\mathcal E^{\le d-1}=\D^{\le d-1}_B$.  The inclusion $\D^{\le d}_B\subseteq\mathcal
E^{\le d}$ is immediate from $S\subseteq B$.  Conversely let $X\in\mathcal E^{\le d}$ and
put
\[
  T=\bigl\{f\in\Lam:\ [X,M_f]\in\mathcal E^{\le d-1}\bigr\}.
\]
Then $T$ is a $\Q(t)$-subspace ($f\mapsto M_f$ is linear); $1\in T$ because
$M_1=\mathrm{id}$ and $[X,\mathrm{id}]=0$; and $S\subseteq T$ by hypothesis.  If
$f,g'\in T$ then $M_{fg'}=M_fM_{g'}$ and
\[
  [X,M_{fg'}]=[X,M_f]\,M_{g'}+M_f\,[X,M_{g'}],
\]
and both summands lie in $\mathcal E^{\le d-1}$ by Step~1; so $fg'\in T$.  Hence $T$ is a
subalgebra containing $1$ and $S$, so $T\supseteq B$.  Thus $[X,M_f]\in\mathcal E^{\le
d-1}=\D^{\le d-1}_B$ for every $f\in B$, i.e.\ $X\in\D^{\le d}_B$.
\end{proof}

\begin{corollary}\label{cor:pe}
$\ord_{\Q(t)[p_e]}(X)\le d$ if and only if $\ad^{\,d+1}(X)=0$, where
\[
  \ad(X):=[X,M_{p_e}],\qquad
  \ad^{\,d}(X)=\sum_{r=0}^{d}(-1)^{r}\binom{d}{r}\,M_{p_e}^{\,r}\,X\,M_{p_e}^{\,d-r},
\]
the second equality because the $M_{p_e}$ commute with one another.  In particular, if
$\ad^{\,d}(X)\neq0$ for every $d\ge1$ then $\ord_{\Q(t)[p_e]}(X)=\infty$, and hence by
Lemma~\ref{lem:mono} also $\ord_{(e)}(X)=\ord(X)=\infty$.
\end{corollary}

\begin{proof}
Lemma~\ref{lem:gen} with $S=\{p_e\}$, $B=\Q(t)[p_e]$.
\end{proof}

\section{The witness}

Fix $e\ge2$ and $d\ge1$.  Write $M_\varnothing=\Z_{<0}$ for the Maya set of the empty
partition, and
\[
  M'_d:=\bigl(\Z_{<0}\setminus\{-2\}\bigr)\cup\bigl\{(d+1)e-2\bigr\}.
\]

\begin{lemma}[the target is a hook]\label{lem:hook}
$M'_d$ is the Maya set of $\mu_d:=\bigl((d+1)e-1,\,1\bigr)$, a partition of $(d+1)e$.
\end{lemma}

\begin{proof}
List the beads of $M'_d$ in decreasing order: $(d+1)e-2$ (which is $\ge e-2\ge0$, so it is
the largest), then $-1,-3,-4,-5,\dots$.  If $\beta_j$ is the $j$-th largest bead then
$\lambda_j=\beta_j+j$.  So $\lambda_1=(d+1)e-2+1=(d+1)e-1$, $\lambda_2=-1+2=1$, and
$\lambda_j=-j+j=0$ for $j\ge3$.  Its size is $(d+1)e$, which is also the degree by which
$\ad^{\,d}(R_e(t))$ raises, since each of the $d+1$ factors adds $e$ cells.
\end{proof}

\begin{lemma}[the channel is a single forced route]\label{lem:unique}
Let $e\ge2$, $d\ge1$.  There is exactly one sequence of $d+1$ legal $e$-moves
\[
  M_\varnothing=M^{(0)}\longrightarrow M^{(1)}\longrightarrow\cdots\longrightarrow M^{(d+1)}=M'_d ,
\]
namely: for $k=0,1,\dots,d$ in this order, move the bead at $-2+ke$ to $-2+(k+1)e$.  The
intermediate states are $M^{(k)}=(\Z_{<0}\setminus\{-2\})\cup\{-2+ke\}$.
\end{lemma}

\begin{proof}
Suppose $x_1\to x_1+e,\ \dots,\ x_{d+1}\to x_{d+1}+e$ is such a sequence.  For $v\in\Z$
let
\[
  \delta(v)=\#\{i:x_i+e=v\}-\#\{i:x_i=v\}.
\]
The occupancy of $v$ increases by $1$ each time a move lands on $v$ and decreases by $1$
each time a move departs from $v$, so summing over the whole sequence,
$\delta(v)=\mathbf 1_{M'_d}(v)-\mathbf 1_{M_\varnothing}(v)$, i.e.
\[
  \delta=\mathbf 1_{\{(d+1)e-2\}}-\mathbf 1_{\{-2\}} .
\]
(This is pure counting; it does not use legality.)

Every move preserves the residue class mod $e$.  Fix a class, identify it with $\Z$ by
$c+ke\leftrightarrow k$, and let $m_k$ be the number of moves $c+ke\to c+(k+1)e$; only
finitely many $m_k$ are nonzero, and $\delta(c+ke)=m_{k-1}-m_k$.

If $c\not\equiv-2\pmod e$ then $\delta$ vanishes on the class, so $m_{k-1}=m_k$ for all
$k$; with finite support this forces $m_k=0$ for all $k$.  So no move occurs outside the
class of $-2$.

Take $c=-2$, indexing so that $k$ corresponds to the site $-2+ke$; then
$\delta(-2+ke)=\mathbf 1_{k=d+1}-\mathbf 1_{k=0}$.  For $k\le-1$ we get $m_{k-1}=m_k$,
hence $m_k=0$ for all $k\le-1$; at $k=0$, $m_{-1}-m_0=-1$ gives $m_0=1$; for $1\le k\le d$,
$m_{k-1}=m_k$ gives $m_k=1$; at $k=d+1$, $m_d-m_{d+1}=1$ gives $m_{d+1}=0$, and then
$m_k=0$ for all $k>d+1$.  So the multiset of moves is exactly
$\{-2+ke\to-2+(k+1)e:0\le k\le d\}$, which already has $d+1$ elements --- the total number
of moves.  Hence there are no others, in any class.

It remains to fix the order.  For $k\ge1$ the site $-2+ke$ satisfies $-2+ke\ge e-2\ge0$,
so it is \emph{unoccupied} in $M_\varnothing$, and the only move that can occupy it is the
move indexed $k-1$.  Legality of move $k$ requires $-2+ke$ to be occupied at that instant,
so move $k$ must come after move $k-1$.  Therefore the order is $0,1,\dots,d$, and the
intermediate states are as stated.
\end{proof}

\begin{lemma}[the two weights]\label{lem:weights}
In the notation of Lemma~\ref{lem:unique}, let $h_k=\#\bigl(M^{(k)}\cap(-2+ke,\,-2+(k+1)e)\bigr)$
be the height picked up by move $k$.  Then, for $e\ge2$,
\[
  h_0=1,\qquad h_k=0\ \ (1\le k\le d).
\]
\end{lemma}

\begin{proof}
$k=0$: $M^{(0)}=\Z_{<0}$ and the open window is $(-2,\,e-2)=\{-1,0,1,\dots,e-3\}$.  Since
$e\ge2$ we have $e-2>-1$, so $-1$ lies in the window; the remaining sites are $\ge0$ and
therefore unoccupied.  Hence $h_0=1$.  (For $e=1$ the window is empty and $h_0=0$; this is
exactly where the argument needs $e\ge2$.)

$1\le k\le d$: the open window is $\{-1+ke,\dots,-3+(k+1)e\}$, and its least element
satisfies $-1+ke\ge-1+e\ge1>0$.  The only non-negative bead of
$M^{(k)}=(\Z_{<0}\setminus\{-2\})\cup\{-2+ke\}$ is $-2+ke$, which is the left endpoint and
so not in the open window.  Hence $h_k=0$.
\end{proof}

\begin{theorem}[main]\label{thm:main}
For every $e\ge2$ and every $d\ge1$, with $\mu_d=\bigl((d+1)e-1,\,1\bigr)$,
\[
  \boxed{\ \bigl\langle s_{\mu_d},\ \ad(M_{p_e})^{\,d}\bigl(R_e(t)\bigr)\,s_\varnothing\bigr\rangle
  \;=\;(-1)^{d}\,(1+t).\ }
\]
The same formula holds with $t$ replaced by any scalar $c$.
\end{theorem}

\begin{proof}
By Corollary~\ref{cor:pe},
$\ad^{\,d}(R_e(t))=\sum_{r=0}^d(-1)^r\binom dr M_{p_e}^{\,r}R_e(t)M_{p_e}^{\,d-r}$.  By
Proposition~\ref{prop:mn}, $M_{p_e}=R_e(-1)$, so every summand is a composite of $d+1$
operators each of which acts on Maya sets by legal $e$-moves
(Proposition~\ref{prop:abacus}); the two possible weights of a move of height $h$ are
$t^{h}$ (if performed by $R_e(t)$) and $(-1)^{h}$ (if performed by $M_{p_e}$).

Fix $r$ and look at the summand $M_{p_e}^{\,r}R_e(t)M_{p_e}^{\,d-r}$.  Its
$\langle M'_d|\cdot|M_\varnothing\rangle$ entry is the sum, over all sequences of $d+1$
legal $e$-moves from $M_\varnothing$ to $M'_d$, of the product of the weights.  By
Lemma~\ref{lem:unique} there is exactly one such sequence, and it is the same one for
every $r$; the operators act rightmost-first, so $R_e(t)$ performs move number
$j:=d-r$ (indices $0,\dots,d$ as in Lemma~\ref{lem:unique}) and $M_{p_e}$ performs all the
others.  By Lemma~\ref{lem:weights} the weight of that unique route is
\[
  w_j=t^{\,h_j}\prod_{i\neq j}(-1)^{h_i}
    =\begin{cases}t^{1}\cdot 1=t, & j=0,\\[2pt] t^{0}\cdot(-1)^{1}=-1, & 1\le j\le d.\end{cases}
\]
Therefore, re-indexing by $j=d-r$,
\[
  \bigl\langle M'_d\bigl|\ad^{\,d}(R_e(t))\bigr|M_\varnothing\bigr\rangle
  =\sum_{j=0}^{d}(-1)^{\,d-j}\binom{d}{j}\,w_j
  =\underbrace{\sum_{j=0}^{d}(-1)^{d-j}\binom dj\,(-1)}_{=-(1-1)^d=0\ (d\ge1)}
   \;+\;(-1)^{d}\binom d0\bigl(t-(-1)\bigr),
\]
where in the first sum we replaced every $w_j$ by $-1$ and in the second we corrected the
$j=0$ term by $w_0-(-1)=t+1$.  This equals $(-1)^d(1+t)$.  Nothing in the argument used
that $t$ was an indeterminate, so it holds verbatim for a scalar $c$.  Finally
$\langle M'_d|\cdot\rangle=\langle s_{\mu_d},\cdot\rangle$ by Lemma~\ref{lem:hook}.
\end{proof}

\begin{corollary}[Q96: the answer is \emph{no}]\label{cor:q96}
For every $e\ge2$ and every scalar $c\neq-1$ --- and for the indeterminate $t$ over $\Q(t)$ ---
\[
  \ord_{\Q(t)[p_e]}\bigl(R_e(c)\bigr)
  =\ord_{(e)}\bigl(R_e(c)\bigr)
  =\ord\bigl(R_e(c)\bigr)=\infty .
\]
At $c=-1$ all three are $0$.
\end{corollary}

\begin{proof}
$(-1)^d(1+c)\neq0$ for every $d$ when $c\neq-1$, so $\ad^{\,d}(R_e(c))\neq0$ for every
$d\ge1$; apply Corollary~\ref{cor:pe} and then Lemma~\ref{lem:mono}.  At $c=-1$,
$R_e(-1)=M_{p_e}$ is a multiplication operator, of order $0$ over any base.
\end{proof}

So shrinking the base from $\Lam$ all the way down to $\Q(t)[p_e]$ --- a single generator,
the smallest base for which the question is not vacuous --- does not make the order finite.
The infinite order of $R_e(t)$ is not an artifact of measuring the deformation in all the
$p_m$ at once.

\begin{remark}[the pin of the earlier proof was removable]\label{rem:pin}
The full-base result \cite[Cor.~4.9]{Clio0906} was proved by iterating $\ad(M_{p_1})$, and
$p_1\notin\Lam^{(e)}$ for $e\ge2$; the entire mechanism of that proof is therefore
unavailable here, which is why Q96 was open rather than a corollary.  The outcome is that
the pinned generator was \emph{not} carrying the result: the same phenomenon is visible
against $M_{p_e}$ alone.  Theorem~\ref{thm:main} moreover recovers
\cite[Cor.~4.9]{Clio0906} in one stroke and uniformly in the parameter --- including the
scalar $c=0$, for which \cite{Clio0906} needed a second theorem
(\cite[Thm.~B]{Clio0906}), because there the leading route weight degenerated.  Here
$c=0$ gives $(-1)^d$, which is visibly nonzero.
\end{remark}

\begin{remark}[why the obvious channel is blind]\label{rem:blind}
It is worth recording the channel that does \emph{not} work, since it is the first one to
try.  Take the target $\bigl((d+1)e\bigr)$, a single row, whose Maya set is
$(\Z_{<0}\setminus\{-1\})\cup\{(d+1)e-1\}$.  The same flow argument shows the route is
again unique --- the bead at $-1$ walks to the right in $d+1$ hops --- but now
\emph{every} window $(-1+ke,\,-1+(k+1)e)$ consists of sites $\ge0$ and is therefore empty,
so all heights are $0$, all weights are $1$, and the matrix element is
$\sum_j(-1)^{d-j}\binom dj=0$ for every $d\ge1$.  This channel is constant in the direction
it is meant to test: it cannot see $t$ at all.  Starting the walk at $-2$ instead of $-1$
is exactly what puts one bead, the one at $-1$, inside the first window.
\end{remark}

\begin{remark}[independence from the LLT map]\label{rem:lam}
Every step above is a statement about legal $e$-moves on Maya sets, the height statistic
$\#(M\cap(b,b+e))$, and a binomial identity.  No LLT function, no $q$-deformation and no
image under Lam's map $\Phi$ appears anywhere.  This matters because $\Phi$ sends the
relevant operators to multiplication operators, which have order $0$; an argument routed
through $\Phi$ could not distinguish order $0$ from order $\infty$ here.  The present
argument does not use it.
\end{remark}

\section{The two-parameter commutator}\label{sec:two}

The $d=0$ instance of Corollary~\ref{cor:q96} is the statement
$[R_e(t),R_e(-1)]\neq0$: one operator, two values of the parameter.  We record the general
two-parameter commutator, which explains the shape of the obstruction and specialises
correctly at $s=t$ and at $t=s=-1$.  The derivation is the route enumeration of
\cite[Thm.~2.3]{Clio0907} carried out with the two parameters kept separate: each route
picks up each height-weight from whichever operator traverses that window.

\begin{theorem}[{matrix elements of $[R_e(t),R_f(s)]$}]\label{thm:two}
Let $e,f\ge1$, let $M$ be a Maya set and $M'\neq M$.  Then
$\langle M'|[R_e(t),R_f(s)]|M\rangle=0$ unless $|M\triangle M'|\in\{2,4\}$, and:

\smallskip
\noindent\textbf{(1) One-bead sector.}  If $M'=M\cdot(a,e+f)$ is a legal $(e{+}f)$-move,
write $N=\#(M\cap(a,a+e+f))$, $A=\#(M\cap(a,a+f))$, $B=\#(M\cap(a,a+e))$, and
$m_e=\mathbf 1[a+e\in M]$, $m_f=\mathbf 1[a+f\in M]$.  Then
\[
  \boxed{\ \langle M'|[R_e(t),R_f(s)]|M\rangle
  = s^{A}t^{\,N-1-A}\bigl(t-m_f(1+t)\bigr)\;-\;t^{B}s^{\,N-1-B}\bigl(s-m_e(1+s)\bigr).\ }
\]

\smallskip
\noindent\textbf{(2) Two-bead sector.}  If $|M\triangle M'|=4$, write
$M\setminus M'=\{b,c\}$.  Call an ordered pair $(b,c)$ (with $\{b,c\}=M\setminus M'$) a
\emph{legal assignment} if $b+e\in M'\setminus M$, $c+f\in M'\setminus M$ and
$b,c,b+e,c+f$ are pairwise distinct.  For such a pair put $P=\#(M\cap(b,b+e))$,
$Q=\#(M\cap(c,c+f))$ and
\[
  k=k^{e,f}_{b,c}=\mathbf 1[\,b+e\in(c,c+f)\,]-\mathbf 1[\,b\in(c,c+f)\,].
\]
Then
\[
  \boxed{\ \langle M'|[R_e(t),R_f(s)]|M\rangle
  \;=\;\sum_{(b,c)\ \mathrm{legal}} t^{\,P-k}\,s^{\,Q}\,\bigl(1-(ts)^{k}\bigr).\ }
\]
There is exactly one legal assignment when $e\neq f$, and exactly two --- $(b,c)$ and
$(c,b)$, whose indices $k$ are negatives of one another --- when $e=f$.
\end{theorem}

\begin{proof}
Both sectors are the enumeration of \cite[Thm.~2.3]{Clio0907}; only the bookkeeping of the
parameters changes, so we give that.

(1) $R_e$ first.  As in \cite[Thm.~2.3, Step~2]{Clio0907}, exactly two route types reach
$M\cdot(a,e+f)$: the $e$-move from $a$ followed by the $f$-move from $a+e$, which exists
iff $m_e=0$ and has heights $B$ (for the $e$-move) and $N-B$ (for the $f$-move, using
$m_e=0$); and the $e$-move from $a+f$ followed by the $f$-move from $a$, which exists iff
$m_f=1$ and has heights $N-1-A$ and $A$.  Attaching $t$ to $e$-moves and $s$ to $f$-moves,
\[
  \langle M'|R_f(s)R_e(t)|M\rangle=(1-m_e)\,t^{B}s^{\,N-B}+m_f\,t^{\,N-1-A}s^{A}.
\]
Exchanging the roles of $(e,t)$ and $(f,s)$ --- the target is unchanged ---
\[
  \langle M'|R_e(t)R_f(s)|M\rangle=(1-m_f)\,s^{A}t^{\,N-A}+m_e\,s^{\,N-1-B}t^{B}.
\]
Subtract and group the two monomial types.  The $s^{A}t^{\bullet}$ terms give
$(1-m_f)s^At^{N-A}-m_f s^At^{N-1-A}=s^At^{N-1-A}\bigl(t-m_f(1+t)\bigr)$, using
$m_f\in\{0,1\}$; the $t^{B}s^{\bullet}$ terms give
$m_es^{N-1-B}t^{B}-(1-m_e)t^Bs^{N-B}=-t^Bs^{N-1-B}\bigl(s-m_e(1+s)\bigr)$.

(2) Fix a legal assignment $(b,c)$: $R_e$ moves $b$ by $e$ and $R_f$ moves $c$ by $f$.  If
the $e$-move goes first its height is $P$ and the subsequent $f$-move sees a Maya set in
which the count on $(c,c+f)$ has changed by exactly $k$, so its height is $Q+k$; the
contribution to $R_f(s)R_e(t)$ is $t^{P}s^{\,Q+k}$.  If the $f$-move goes first its height
is $Q$ and the subsequent $e$-move has height $P+(\beta-\alpha)$ where, by
\cite[Lem.~2.1]{Clio0907}, $\beta-\alpha=-k$ because $b\neq c$ and $b+e\neq c+f$; the
contribution to $R_e(t)R_f(s)$ is $s^{Q}t^{\,P-k}$.  The net contribution to the
commutator is $s^{Q}t^{P-k}-t^{P}s^{Q+k}=t^{P-k}s^{Q}(1-(ts)^{k})$.  Summing over legal
assignments gives the formula.

For the count of legal assignments: the bead moved by $R_e$ is displaced by $e$ and the one
moved by $R_f$ by $f$.  If $e\neq f$ this determines which is which, since the alternative
would force $c-b=f-e$ and $c-b=e-f$ simultaneously.  If $e=f$ both assignments are
available, and interchanging $b$ and $c$ replaces $k^{e,e}_{b,c}$ by $k^{e,e}_{c,b}=-k$.
\end{proof}

\begin{corollary}[the specialisations]\label{cor:spec}
\leavevmode
\begin{enumerate}
\item[(i)] At $s=t$ the formulas reduce to \cite[Thm.~2.3]{Clio0907}: sector (1) becomes
$(m_e-m_f)(1+t)t^{N-1}$ and sector (2) becomes $t^{P+Q}(t^{-k}-t^{k})$ summed over legal
assignments, which cancels in pairs when $e=f$, consistently with $[R_e(t),R_e(t)]=0$.
\item[(ii)] At $t=s=-1$ both sectors vanish identically, consistently with the fact that
all $M_{p_k}$ commute.
\item[(iii)] The two-bead sector vanishes identically \emph{if and only if} $ts=1$.
\item[(iv)] $[R_e(t),R_e(s)]$ has a nonzero two-bead part whenever $t\neq s$ and
$ts\neq1$.
\end{enumerate}
\end{corollary}

\begin{proof}
(i) In sector (1) put $s=t$: the first term is $t^{N-1}(t-m_f(1+t))$ and the second is
$t^{N-1}(t-m_e(1+t))$, whose difference is $t^{N-1}(1+t)(m_e-m_f)$.  In sector (2),
$t^{P-k}t^Q(1-t^{2k})=t^{P+Q}(t^{-k}-t^{k})$; at $e=f$ the two legal assignments
contribute $t^{P+Q}(t^{-k}-t^k)$ and $t^{P+Q}(t^{k}-t^{-k})$.

(ii) At $t=-1$ we have $t-m_f(1+t)=-1$ and likewise $s-m_e(1+s)=-1$, so sector (1) becomes
$(-1)^{A}(-1)^{N-1-A}(-1)-(-1)^{B}(-1)^{N-1-B}(-1)=0$; and $ts=1$ kills sector (2) by (iii).

(iii) Each summand carries the factor $1-(ts)^{k}$, which vanishes when $ts=1$.
Conversely $k=\pm1$ does occur (for instance $e=1$, $f=2$ and a configuration with
$b\in(c,c+f)$), and there $1-(ts)^{\pm1}\neq0$ unless $ts=1$.

(iv) Immediate from (iii) and the two-assignment description: at $e=f$ the two summands are
$t^{P-k}s^{Q}(1-(ts)^{k})$ and $t^{Q+k}s^{P}(1-(ts)^{-k})$, whose sum is
$t^{P-k}s^{Q}-t^{P}s^{Q+k}+t^{Q+k}s^{P}-t^{Q}s^{P-k}$; this is not identically zero.
\end{proof}

\begin{remark}[a correction to a case split]\label{rem:casesplit}
Part (iv) corrects a reading of \cite[Thm.~2.3(2)]{Clio0907}, whose clause ``if $e=f$ then
the matrix element is $0$'' is a \emph{one-parameter} statement: it holds because the two
legal assignments cancel when $s=t$, not because the sector is empty.  Once $t\neq s$ the
cancellation fails.  Concretely, for $e=f=3$, $\lambda=\varnothing$, $b=6$, $c=8$ in a
$9$-bead window one finds $P=2$, $Q=0$, $k=1$ and matrix element
$s^{2}t-st^{2}-s+t=(t-s)(1-st)$, which is nonzero off the two loci $s=t$ and $st=1$.  This
was found by the numerical sweep of \S\ref{sec:verif}, which reported exactly $61$
disagreements against a first version of the formula that had imported the $e=f$ clause
verbatim; all $61$ were $e=f$ two-bead entries.
\end{remark}

\begin{remark}[why the one-bead sector no longer factors]\label{rem:twocolour}
In operator form the one-bead sector of $[R_e(t),R_f(s)]$ is
\[
  \sum_{a\in\Z}\psi_{a+e+f}\psi_a^{*}\Bigl[
     (-s)^{N_{(a,a+f)}}(-t)^{N_{(a+f,\,a+e+f)}}
   \;-\;(-t)^{N_{(a,a+e)}}(-s)^{N_{(a+e,\,a+e+f)}}\Bigr],
\]
where $N_{(x,y)}=\sum_{x<j<y}n_j$.  Each term dresses the interval $(a,a+e+f)$ in
\emph{two colours}: the sub-interval traversed by $R_e$ carries $t$, the one traversed by
$R_f$ carries $s$ --- and the two orderings split the interval at different points
($a+f$ versus $a+e$).  When $s=t$ the colours merge and both terms collapse to
$(-t)^{N_{(a,a+e+f)}}$ times a local factor, which is how the one-parameter answer acquires
its clean form $-\frac{1+t}{t}\Phi_{e,f}$.  When $s\neq t$ no such collapse occurs.  This
is the structural reason a change of parameter cannot be gauged away, and it is the
mechanism behind Theorem~\ref{thm:main}: in the witness channel the first window
(the one containing the bead $-1$) is traversed by $R_e(t)$ in exactly one of the $d+1$
routes, and by $M_{p_e}$ in all the others.
\end{remark}

\section{Verification}\label{sec:verif}

All computations were run with two code-disjoint engines and, for
Proposition~\ref{prop:mn}, a third.

\begin{enumerate}
\item \textbf{Engine A} (abacus) versus \textbf{engine B} (brute-force enumeration of
border strips on Young diagrams, with explicit $2\times2$ and connectivity tests):
$R_g(t)s_\lambda$ agrees on $60/60$ pairs, $1\le g\le5$, $\lambda$ ranging over $12$
partitions with $|\lambda|\le8$, with $t$ an indeterminate.
\item \textbf{Premise} $M_{p_g}=R_g(-1)$ (Proposition~\ref{prop:mn}), re-verified against
\textbf{engine C}: $s_\lambda$ built by Jacobi--Trudi as a polynomial in $p_1,p_2,\dots$,
multiplied by the \emph{symbol} $p_g$, and re-expanded in the Schur basis by solving a
linear system over the $p$-monomial basis.  Agreement on $29/29$ cases, $1\le g\le4$ and
$|\lambda|+g\le7$.  Engine C shares no code with A or B.
\item \textbf{Theorem~\ref{thm:main}}: the predicted value $(-1)^d(1+t)$ was written down
from Lemmas~\ref{lem:unique}--\ref{lem:weights} \emph{before} being computed, and then
confirmed on $13/13$ cases: $(e,d)$ with $2\le e\le5$, $1\le d\le5$ and $(d+1)e\le15$.
\item \textbf{Theorem~\ref{thm:main}, independent path}: the same $5$ cases
($e=2,d\le3$ and $e=3,d\le2$) recomputed with engines B and C only --- border strips on
Young diagrams for $R_e(t)$, literal multiplication by $p_e$ in $\Q(t)[p_1,p_2,\dots]$ for
$M_{p_e}$ --- with no abacus anywhere in the path.  $5/5$.
\item \textbf{Scalar specialisations}: $c\in\{0,1,2,-\tfrac12\}$ give $(-1)^d(1+c)$ on all
$6$ available $(e,d)$; $c=-1$ gives $0$ on all $6$.
\item \textbf{Theorem~\ref{thm:two}}: predicted versus computed matrix elements of
$[R_e(t),R_f(s)]$, agreement on $\mathbf{1829/1829}$ entries.  The denominator factors as
$1232$ one-bead entries $+\ 597$ two-bead entries, of which $61$ have $e=f$; those $61$ are
exactly the entries that a version of the formula importing the one-parameter $e=f$ clause
got wrong (Remark~\ref{rem:casesplit}).  Quantifier: $1\le e,f\le4$; all $30$ partitions
with $|\lambda|\le6$; every nonzero target of $[R_e(t),R_f(s)]s_\lambda$.
\item \textbf{Corollary~\ref{cor:spec}}: specialising $s\to t$ reproduces an independently
computed one-parameter commutator on $304/304$ source states ($1\le e,f\le4$,
$|\lambda|\le5$); $t=s=-1$ annihilates all $304$; $s=1/t$ annihilates all $379$ two-bead
entries.
\end{enumerate}

\subsection*{Negative controls}

Both controls move, i.e.\ each returns a different answer from the one the theorem returns.

\begin{itemize}
\item \textbf{$e=1$.}  Here $\Lam^{(1)}=\Lam$ and the theorem must \emph{not} apply.  It
does not: a border strip of size $1$ is a single cell, of height $0$, so $R_1(t)=R_1(-1)=
M_{p_1}$ is independent of $t$ and has order $0$ over every base.  Checked: every
coefficient of $R_1(t)s_\lambda$ equals $1$ for all $|\lambda|\le7$; $[R_1(t),M_{p_m}]=0$
for $1\le m\le4$ and all $|\lambda|\le5$; and the same machinery that produces
$(-1)^d(1+t)$ for $e\ge2$ returns $\ad^{\,d}=0$ at $e=1$.  The proof localises the failure
exactly: Lemma~\ref{lem:weights} needs $e\ge2$ for $-1$ to lie in the first window.
\item \textbf{$t=-1$.}  $\ad(M_{p_e})^{d}(R_e(-1))=0$ for $e\in\{2,3\}$, $1\le d\le3$ and
all $|\lambda|\le4$, against $(-1)^d(1+t)\neq0$ for generic $t$.  The machinery therefore
distinguishes the anchor from its neighbourhood, and is not constant in the direction it is
meant to test.
\end{itemize}

\begin{remark}[an instrument failure, recorded]
The first run of verification item~4 disagreed with item~3, returning $-2$ where the
theorem predicts $-t-1$.  The disagreement was in the instrument: the routine extracting
$p$-monomial coefficients used \texttt{as\_coefficients\_dict}, which returns only the
\emph{numeric} part of a coefficient, silently discarding the symbolic factor $t$ --- so it
was reporting the value at $t=1$, and indeed $-t-1$ at $t=1$ is $-2$.  After the extractor
was rewritten to make the $p_k$ the polynomial generators and keep the full coefficient,
all $5$ cases agreed.  Verification items~1 and~2 were unaffected, since they involve no
symbolic coefficients ($t=-1$ there).
\end{remark}

\section{What is settled and what is not}\label{sec:gaps}

\textbf{Settled.}  Q96 is answered, in the negative and in the strong form: the order is
infinite already over $\Q(t)[p_e]$ (Corollary~\ref{cor:q96}).  The reduction to algebra
generators is proved, not imported (Lemma~\ref{lem:gen}).  The two-parameter commutator is
in closed form (Theorem~\ref{thm:two}) with its specialisations checked
(Corollary~\ref{cor:spec}).

\textbf{Consequence for the programme.}  The reading of Q96 under which a ``yes'' would
have given Q91 a conditional normal form over $\Lam^{(e)}$ is closed off: the obstruction
is not removable by restricting the base.  Whatever normal form $R_e(t)$ has, it is not
one with finitely many $\partial_{p_{me}}$.

\textbf{Not settled, and not attempted here.}
\begin{enumerate}
\item $\ord_{(e)}\bigl(R_g(t)\bigr)$ for $g$ not a multiple of $e$, and more generally the
whole two-index table.  The flow argument of Lemma~\ref{lem:unique} uses that every move
has the \emph{same} size $e$; with two step sizes the multiset of moves is no longer forced
by divergence alone and the route count is genuinely harder.
\item The locus $ts=1$ of Corollary~\ref{cor:spec}(iii).  It is a second distinguished
locus for the two-parameter family, alongside $s=t$ and the anchor $t=s=-1$ (which lies on
both).  We have proved only that the two-bead sector dies there; what the surviving
one-bead operator is, and whether the locus means anything, is not examined.
\item Nothing in this note is a claim about the literature.  No source was read during this
session beyond the author's own earlier files, and no priority or novelty assertion is made
about any statement above.
\end{enumerate}

\begin{thebibliography}{9}
\bibitem{Clio0903} Clio, \emph{The multiplication defect of the ribbon operator} (Q75),
\texttt{proofs/2026-09-03-Q75-multiplication-defect.tex}, 3 September 2026.
\bibitem{Clio0906} Clio, \emph{$\ord(N_e)=\infty$, and the order filtration is
discontinuous at $t=-1$} (Q84), \texttt{proofs/2026-09-06-Q84-order-of-N-e.tex},
6 September 2026.
\bibitem{Clio0907} Clio, \emph{The cross-rank commutator $[R_e(t),R_f(t)]$ in closed form}
(Q92), \texttt{proofs/2026-09-07-Q92-cross-rank-commutator.tex}, 7 September 2026.
\end{thebibliography}

\end{document}
